\section{Analysis of Workflow Nets}

\emph{Soundness} characterises a well-formed business process, and the Petri-net machinery developed so far studies it.\footnote{Weske M., \emph{Business Process Management}, Springer, 2007, Ch.~6.} Boundedness and liveness suit an \emph{ongoing} system; a workflow instead has a definite start and end per case, and soundness adapts them to that setting. Since BPMN, EPC and workflow nets share the same building blocks (start/end markers, splits/joins, tasks, links), the analysis is developed on workflow nets and transfers to the other notations through the usual mappings; the reseller process of the orchestration chapter (drawn in all three notations) is the running motivation.

\subsection{Workflow nets, structurally}

Recall the workflow-net definition from the Petri-nets-basics chapter: a Petri net with a distinguished initial place $i$ (${}^\bullet i = \emptyset$), a distinguished final place $o$ ($o^\bullet = \emptyset$), and every other node on a path from $i$ to $o$. Places hold tokens representing \emph{cases}, transitions are \emph{activities}, a doubled border marks an activity refined by another workflow net, and the canonical initial marking is one token on $i$.

The definition is purely structural, yet condition 3 already rejects several malformed shapes: a net with \emph{no} source or sink, a net with \emph{several} sources or sinks, and \emph{dangling tasks} without inputs or outputs. What it does not catch are vacuous decorations (a ``split'' with one output, a ``join'' with one input): syntactically legal, semantically misleading, left to later analysis. Terminology: a \emph{net system} is a net with an initial marking, the setting of behavioural properties; a \emph{net} alone is the setting of structural ones. For WF nets the two nearly coincide thanks to the canonical marking.

\subsection{Behavioural pathologies}

Structural correctness does not imply behavioural correctness. Four small, structurally impeccable workflow nets exhibit four distinct failure modes:
\begin{itemize}
  \item \textbf{dead transition}: a synchronising transition whose input places can never be simultaneously marked (e.g.\ an AND-join downstream of an XOR-split) never fires: a path no case can travel;
  \item \textbf{pending tokens upon completion}: a run marks the final place $o$ while leaving tokens behind on intermediate places: the case is reported complete, but the net is polluted for the next case;
  \item \textbf{pending activities upon completion}: a back-loop keeps upstream transitions enabled after $o$ is marked: the case is closed, yet work continues;
  \item \textbf{deadlock}: a reachable marking enables nothing and $o$ is unmarked: the case is irrevocably stuck. A larger net can ``look fine'' and still hide one: only replaying the dynamics (rolling back in time through the firing sequence) exposes the marking where the mismatched split/join strands its tokens.
\end{itemize}

\begin{definitionbox}{Livelock}
A \emph{livelock} traps a case in a cycle with no opportunity to end: from some reachable marking, every firing sequence keeps moving tokens around a sub-region without ever marking $o$. The system is not deadlocked (work happens) but the case progresses in circles: a failure class distinct from dead transitions and residual tokens.
\end{definitionbox}

All these flaws are detectable \emph{without knowing the goal} of the business process: they are intrinsic to the net's shape and dynamics, hence amenable to automatic model-level checking. This is \textbf{verification} in the sense of the BPM lifecycle (qualitative questions about the model: can it deadlock? do all cases terminate? can this task ever run?), as opposed to validation against reality, which is the business of process mining. The mismatched splits and joins behind these defects could often be prevented at design time; the structural conditions that guarantee prevention (free-choiceness, S/T-decompositions) come in the following chapters.

\subsection{The language of a workflow net}

\begin{definitionbox}{Language of a WF net}
For a workflow net $N$ with initial place $i$ and final place $o$,
\[
  L(N) = \{\, \sigma \mid i \xrightarrow{\sigma} o \,\}
\]
is the set of firing sequences leading from the marking $i$ to the marking $o$: the \emph{admissible traces}, the legitimate complete runs of a case.
\end{definitionbox}

\begin{examplebox}{Question time: languages of small nets}
\begin{center}
\begin{tikzpicture}[node distance=0.55cm and 0.55cm]
  \node[bpmn event] (i) {}; \fill[accent] (i.center) circle (1.6pt);
  \node[bpmn task, rounded corners=0pt, minimum width=0.4cm, right=of i] (A) {\tiny $A$};
  \node[bpmn event, right=of A] (p) {};
  \node[bpmn task, rounded corners=0pt, minimum width=0.4cm, above right=0.3cm and 0.5cm of p] (B) {\tiny $B$};
  \node[bpmn event, right=of B] (q) {};
  \node[bpmn task, rounded corners=0pt, minimum width=0.4cm, right=of q] (C) {\tiny $C$};
  \node[bpmn task, rounded corners=0pt, minimum width=0.4cm, below right=0.3cm and 1.5cm of p] (D) {\tiny $D$};
  \node[bpmn event, above right=0.3cm and 0.5cm of D] (r) {};
  \node[bpmn task, rounded corners=0pt, minimum width=0.4cm, right=of r] (E) {\tiny $E$};
  \node[bpmn event, right=of E] (o) {};
  \draw[bpmn flow] (i) -- (A); \draw[bpmn flow] (A) -- (p);
  \draw[bpmn flow] (p) -- (B); \draw[bpmn flow] (B) -- (q); \draw[bpmn flow] (q) -- (C);
  \draw[bpmn flow] (p) -- (D);
  \draw[bpmn flow] (C) -- (r); \draw[bpmn flow] (D) -- (r);
  \draw[bpmn flow] (r) -- (E); \draw[bpmn flow] (E) -- (o);
  \begin{scope}[xshift=8.2cm]
  \node[bpmn event] (i2) {}; \fill[accent] (i2.center) circle (1.6pt);
  \node[bpmn task, rounded corners=0pt, minimum width=0.4cm, right=of i2] (A2) {\tiny $A$};
  \node[bpmn event, right=of A2] (p2) {};
  \node[bpmn task, rounded corners=0pt, minimum width=0.4cm, right=of p2] (B2) {\tiny $B$};
  \node[bpmn event, above right=0.3cm and 0.4cm of B2] (q2a) {};
  \node[bpmn event, below right=0.3cm and 0.4cm of B2] (q2b) {};
  \node[bpmn task, rounded corners=0pt, minimum width=0.4cm, right=of q2a] (C2) {\tiny $C$};
  \node[bpmn task, rounded corners=0pt, minimum width=0.4cm, right=of q2b] (D2) {\tiny $D$};
  \node[bpmn event, below right=0.3cm and 0.4cm of C2] (r2a) {};
  \node[bpmn event, above right=0.3cm and 0.4cm of D2] (r2b) {};
  \node[bpmn task, rounded corners=0pt, minimum width=0.4cm] (E2) at ($(r2a)!0.5!(r2b)+(1.0,0)$) {\tiny $E$};
  \node[bpmn event, right=of E2] (o2) {};
  \draw[bpmn flow] (i2) -- (A2); \draw[bpmn flow] (A2) -- (p2); \draw[bpmn flow] (p2) -- (B2);
  \draw[bpmn flow] (B2) -- (q2a); \draw[bpmn flow] (B2) -- (q2b);
  \draw[bpmn flow] (q2a) -- (C2); \draw[bpmn flow] (q2b) -- (D2);
  \draw[bpmn flow] (C2) -- (r2a); \draw[bpmn flow] (D2) -- (r2b);
  \draw[bpmn flow] (r2a) -- (E2); \draw[bpmn flow] (r2b) -- (E2);
  \draw[bpmn flow] (E2) -- (o2);
  \end{scope}
\end{tikzpicture}
\end{center}
Left, a choice: after $A$, either the long path $B\,C$ or the shortcut $D$, then $E$: $L(N) = \{ABCE,\, ADE\}$. Right, a parallel split: $B$ forks $C$ and $D$, joined by $E$; the branches interleave freely: $L(N) = \{ABCDE,\, ABDCE\}$. A loop produces a Kleene-star factor:
\begin{center}
\begin{tikzpicture}[node distance=0.55cm and 0.55cm]
  \node[bpmn event] (i) {}; \fill[accent] (i.center) circle (1.6pt);
  \node[bpmn task, rounded corners=0pt, minimum width=0.4cm, right=of i] (A) {\tiny $A$};
  \node[bpmn event, right=of A] (p) {};
  \node[bpmn task, rounded corners=0pt, minimum width=0.4cm, right=of p] (B) {\tiny $B$};
  \node[bpmn event, right=of B] (q) {};
  \node[bpmn task, rounded corners=0pt, minimum width=0.4cm, right=of q] (D) {\tiny $D$};
  \node[bpmn event, right=of D] (r) {};
  \node[bpmn task, rounded corners=0pt, minimum width=0.4cm, right=of r] (E) {\tiny $E$};
  \node[bpmn event, right=of E] (o) {};
  \node[bpmn task, rounded corners=0pt, minimum width=0.4cm, below=0.5cm of D] (C) {\tiny $C$};
  \draw[bpmn flow] (i) -- (A); \draw[bpmn flow] (A) -- (p); \draw[bpmn flow] (p) -- (B);
  \draw[bpmn flow] (B) -- (q); \draw[bpmn flow] (q) -- (D); \draw[bpmn flow] (D) -- (r);
  \draw[bpmn flow] (r) -- (E); \draw[bpmn flow] (E) -- (o);
  \draw[bpmn flow] (r.south) to[bend left=20] (C.east);
  \draw[bpmn flow] (C.west) to[bend left=20] (p.south);
\end{tikzpicture}
\end{center}
$L(N) = \{ABD(CBD)^k E \mid k \geq 0\}$. Further variants raise three separate points.

A structurally valid net can have $L(N) = \emptyset$: if its only synchronising join requires two tokens that no run can provide simultaneously, \emph{zero} traces complete, and the workflow admits no legitimate case at all.

Matching a regular expression against $L(N)$ needs care with the boundary cases. For a net that runs $A$, then $B$, then either exits with $C$ or iterates $R_1 R_2$ back into $B$, the correct descriptions are $A\,B\,(R_1 R_2 B)^k\, C$ and the equivalent $A\,(B\, R_1 R_2)^k\, B\, C$ ($k \geq 0$); the tempting $A\,(B\, R_1 R_2)^k\, C$ is wrong because $k = 0$ would skip the mandatory $B$.

Counting questions (how often can a task fire? can two occurrences be adjacent? can $B$ outnumber $A$?) are answered by token-flow conservation (each $B$ consumes a token some $A$ produced, so $B$s never outnumber $A$s), which the invariants of the next chapter make precise.
\end{examplebox}

\subsection{Soundness}

\begin{definitionbox}{Soundness}
Informally, a business process is \emph{sound} if it contains no unnecessary tasks, every started case can always run to completion, and nothing is left pending when it completes. Formally, a workflow net $N$ with initial place $i$ and final place $o$ is \textbf{sound} if:
\begin{enumerate}
  \item \textbf{no dead task}: $\forall t \in T,\ \exists M \in [i\rangle :\ M \xrightarrow{t}$;
  \item \textbf{option to complete}: $\forall M \in [i\rangle,\ \exists M' \in [M\rangle :\ M'(o) \geq 1$;
  \item \textbf{proper completion}: $\forall M \in [i\rangle :\ M(o) \geq 1 \Rightarrow M = o$.
\end{enumerate}
Here $i$ and $o$ also denote the markings with a single token on the respective place. The three conditions kill the whole pathology catalogue: dead transitions (1), deadlocks and livelocks (2), residual tokens and activities (3).
\end{definitionbox}

Condition 2 is the subtle one: it does not force every run to terminate, only that completion remains \emph{possible} from every reachable marking: loops are fine. To make this operational on looping nets one adopts the \textbf{fairness assumption}: a transition enabled infinitely often eventually fires, so no task is postponed forever and every fair run of a sound net eventually takes the exit.

\emph{Simulation} (playing the token game, recording runs, randomising choices, all supported by WoPeD) finds defects but cannot prove their absence. \emph{Reachability analysis} is the exhaustive answer: all runs live in the reachability graph, so end states and dangerous states are searchable. Its limit is \emph{state explosion}: the graph can be exponentially larger than the net.

Checking the definition on that graph is expensive: condition 1 is one reachability question \emph{per transition}; condition 2 quantifies over \emph{all reachable markings}; condition 3 is a \emph{non}-reachability claim over the infinitely many $M > o$. The brute-force procedure checks the net is a WF net, builds the RG from $i$, and verifies that every transition labels some arc (condition 1) and that the RG has exactly one sink, equal to $o$ and reachable from every node, with no other node marking $o$ (conditions 2--3). Automatic, but limited by state explosion, useless for \emph{repairing} an unsound model, and inapplicable when the RG is infinite. Hence the structural strategy: translate soundness into liveness and boundedness, which have well-studied algorithms.

\subsection{The auxiliary net $N^*$}

A sound workflow plays a case from $i$ to $o$; if it could be \emph{reset} (a token reaching $o$ moved back to $i$) it could play again, any number of times, turning the open workflow into a cyclic system:
\begin{center}
\begin{tikzpicture}[node distance=0.8cm and 1.0cm]
  \node[bpmn event, label=below:{$i$}] (i) {};
  \fill[accent] (i.center) circle (2pt);
  \node[bpmn task, rounded corners=0pt, right=of i, minimum width=0.5cm] (t1) {};
  \node[bpmn task, rounded corners=0pt, right=of t1, minimum width=1.8cm, minimum height=1.2cm] (bp) {Business Process};
  \node[bpmn task, rounded corners=0pt, right=of bp, minimum width=0.5cm] (t2) {};
  \node[bpmn event, right=of t2, label=below:{$o$}] (o) {};
  \node[bpmn task, rounded corners=0pt, above=1.1cm of bp, minimum width=0.6cm] (rst) {reset};
  \draw[bpmn flow] (i) -- (t1);
  \draw[bpmn flow] (t1) -- (bp);
  \draw[bpmn flow] (bp) -- (t2);
  \draw[bpmn flow] (t2) -- (o);
  \draw[bpmn flow] (o.north) to[bend right=25] (rst.east);
  \draw[bpmn flow] (rst.west) to[bend right=25] (i.north);
\end{tikzpicture}
\end{center}

\begin{definition}[Closure $N^*$]
For a workflow net $N : i \to o$, let $N^*$ be $N$ plus one new transition $\text{reset}$ with ${}^\bullet\text{reset} = \{o\}$ and $\text{reset}^\bullet = \{i\}$.
\end{definition}

\begin{proposition}[$N^*$ is strongly connected]
$N^*$ is strongly connected.
\end{proposition}

\textit{Proof.} Take nodes $x, y$ of $N^*$; we exhibit a directed path from $y$ to $x$. If neither is reset: by WF-net condition 3, $y$ lies on a path $i \to^\ast y \to^\ast o$ and $x$ on a path $i \to^\ast x \to^\ast o$; compose $y \to^\ast o \xrightarrow{\text{reset}} i \to^\ast x$. If $x = \text{reset}$: any path $y \to^\ast o$ extends by the arc $o \to \text{reset}$. If $y = \text{reset}$: the arc $\text{reset} \to i$ extends by a path $i \to^\ast x$. \hfill$\square$

\begin{theorem}[Main theorem, van der Aalst]
$N$ is \textbf{sound} $\iff$ $N^*$ is \textbf{live and bounded}.
\end{theorem}

\textit{Proof ($\Leftarrow$): $N^*$ live and bounded $\Rightarrow$ $N$ sound.} Liveness of $N^*$ gives, for each $t \in T$, a reachable marking enabling $t$: condition 1. Now take any $M \in [i\rangle$ enabling reset; then $M \supseteq o$, and firing reset gives $M' \supseteq i$ with $M' \in [i\rangle$. If $M' \neq i$, the surplus $M' - i$ could be pumped by repeating the run-and-reset cycle, contradicting boundedness of $N^*$; hence $M' = i$ and $M = o$. So whenever $o$ is marked the rest of the net is empty (condition 3), and liveness of reset means a marking enabling it (i.e.\ reaching $o$) is reachable from every reachable marking (condition 2). \hfill$\square$

The other direction needs an auxiliary result.

\begin{lemma}[Technical lemma]
If $N$ is sound, then $M$ is reachable in $N$ iff $M$ is reachable in $N^*$.
\end{lemma}

\textit{Proof.} ($\Rightarrow$) trivial: $N^*$ contains $N$. ($\Leftarrow$) Let $i \xrightarrow{\sigma} M$ in $N^*$; induct on the number $r$ of resets in $\sigma$. If $r = 0$, $\sigma$ runs in $N$. If $r > 0$, split $\sigma = \sigma'\, \text{reset}\, \sigma''$ at the first reset: $\sigma'$ is reset-free, so it runs in $N$ and, since reset was enabled after it, reaches a marking $\supseteq o$; by soundness (proper completion) that marking \emph{is} $o$, and firing reset returns exactly to $i$. Then $i \xrightarrow{\sigma''} M$ in $N^*$ with $r - 1$ resets, and the inductive hypothesis closes. \hfill$\square$

Modulo soundness, reset adds no new states.

\textit{Proof ($\Rightarrow$): $N$ sound $\Rightarrow$ $N^*$ bounded and live.} \emph{Bounded.} Suppose not: some $i \to^\ast M \to^\ast M'$ in $N^*$ with $M \subset M'$; let $L = M' - M > \mathbf{0}$. By the technical lemma $M, M'$ are reachable in $N$; by option to complete, $M \xrightarrow{\sigma} o$ in $N$ for some $\sigma$; by monotonicity, $M' = M + L \xrightarrow{\sigma} o + L$. So $o + L$ is reachable, marks $o$, and differs from $o$: contradicting proper completion.
\par\smallskip
\emph{Live.} Take any transition $t$ of $N^*$ and any reachable $M$; by the technical lemma $M$ is reachable in $N$. Option to complete gives $M \xrightarrow{\sigma} o$; no-dead-task gives $i \xrightarrow{\sigma'} M'' \xrightarrow{t}$ for some $M''$. Then in $N^*$: $M \xrightarrow{\sigma} o \xrightarrow{\text{reset}} i \xrightarrow{\sigma'} M'' \xrightarrow{t}$: from any reachable marking, $t$ can be re-enabled. (For $t = \text{reset}$, stop at $o$.) \hfill$\square$

\subsection{Strong connectedness}

\begin{definitionbox}{Paths, circuits, connectedness}
In a net $(P, T, F)$, a \emph{path} is a non-empty node sequence $x_1 \cdots x_k$ with $(x_i, x_{i+1}) \in F$; a \emph{circuit} is a path with distinct nodes and $(x_k, x_1) \in F$: at least two nodes, since the flow relation has no self-arcs. An \emph{undirected path} allows steps in $F \cup F^{-1}$. A net is \emph{weakly connected} if any two nodes are joined by an undirected path (equivalently: it does not split into separated components), and \emph{strongly connected} if any two are joined by a directed path (equivalently, for weakly connected nets: every arc lies on a directed cycle). Strong implies weak; the converse fails, and ``strongly but not weakly'' is vacuous.
\end{definitionbox}

The producer--consumer family illustrates the cases: with a one-way buffer the net is weakly but not strongly connected (nothing flows back); with producer and consumer disjoint it is not even weakly connected; closing the loop with a return arc makes it strongly connected. Disconnected nets decompose, so weak connectedness is assumed implicitly from now on.

\begin{theorem}[Strong connectedness]
If a weakly connected system is \emph{live and bounded}, then it is \emph{strongly connected}. (Proof omitted.)
\end{theorem}

The contrapositive is a cheap structural filter: a weakly-but-not-strongly connected net cannot be both live and bounded: if live then unbounded, if bounded then not live, and possibly neither (all three combinations occur on small examples). This is also why the $N^*$ construction makes sense: $N$ itself is never strongly connected ($i$ has no inputs), so it can never be live and bounded: only its closure $N^*$ can.

\begin{notebox}{Tool support: WoPeD diagnostics}
WoPeD\footnote{\url{http://woped.dhbw-karlsruhe.de/woped/}} implements the soundness check through the $N^*$ construction: it verifies the workflow-net property, builds $N^*$ internally, and checks boundedness and liveness. \textbf{Reading caveat}: the boundedness/liveness entries of the analysis pane refer to $N^*$, not to the open net $N$: a transition flagged non-live in $N^*$ signals a soundness violation of $N$. The diagnostics localise failures: a \emph{possible deadlock} lists the non-live transitions and highlights the join whose inputs cannot be marked together; an \emph{unbounded place} (``boundedness fails because the case can end leaving a token in $c_8$'') pinpoints residual tokens after completion. On sound nets the pane is all green (zero unbounded places, zero dead or non-live transitions) as on the claim-handling workflow (register, classify, simple/complex, parallel insurance and garage checks, decide, pay/send letter), which additionally passes the structural battery (free-choiceness, S-components, well-structuredness) of the next chapters.
\end{notebox}

\subsection{The Computer Repair Service}

\begin{examplebox}{CRS workflow: from text to a sound net}
\emph{Specification}: a customer brings a defective computer; the CRS checks the defect and produces a repair-cost estimate. If the cost is unacceptable the customer takes the computer home unrepaired; otherwise the repair proceeds with two activities in arbitrary order (check/repair the hardware, check/configure the software) after which the system is tested: on error the repair repeats, otherwise the computer is returned.
\begin{center}
% TODO: aesthetic pass: the net is very wide and squeezed by resizebox;
% consider a two-row layout or smaller node spacing.
\resizebox{\linewidth}{!}{%
\begin{tikzpicture}
  \node[bpmn event, label=below:{\scriptsize $i$}] (i) at (0,0) {};
  \fill[accent] (i.center) circle (1.6pt);
  \node[bpmn task, rounded corners=0pt, font=\tiny\sffamily, align=center] (rc) at (1.4,0) {receive\\ computer};
  \node[bpmn event] (p1) at (2.7,0) {};
  \node[bpmn task, rounded corners=0pt, font=\tiny\sffamily, align=center] (cd) at (3.9,0) {check\\ defects};
  \node[bpmn event] (p2) at (5.1,0) {};
  \node[bpmn task, rounded corners=0pt, font=\tiny\sffamily, align=center] (ec) at (6.3,0) {estimate\\ cost};
  \node[bpmn event, label=below:{\scriptsize $p_3$}] (p3) at (7.5,0) {};
  \node[bpmn task, rounded corners=0pt, font=\tiny\sffamily, align=center] (te) at (8.9,-1.3) {too\\ expensive};
  \node[bpmn task, rounded corners=0pt, font=\tiny\sffamily, align=center] (ac) at (8.9,0.9) {acceptable\\ cost};
  \node[bpmn event, label=above:{\scriptsize $p_5$}] (p5) at (10.3,0.9) {};
  \node[bpmn task, rounded corners=0pt, font=\tiny\sffamily, align=center] (sr) at (11.6,0.9) {start\\ repair};
  \node[bpmn event] (p6) at (12.9,1.8) {};
  \node[bpmn task, rounded corners=0pt, font=\tiny\sffamily, align=center] (chw) at (14.1,1.8) {check\\ hw};
  \node[bpmn event, label=above:{\scriptsize $p_7$}] (p7) at (15.3,1.8) {};
  \node[bpmn task, rounded corners=0pt, font=\tiny\sffamily, align=center] (rhw) at (16.5,1.8) {repair\\ hw};
  % Fixed 2026-07-09: TikZ node names were off-by-one from their visible labels (node (p10) carried label p_{11}, (p11) carried p_{12}); renamed nodes to (p11n)/(p12n) to match, so the source is not misleading. Labels and topology unchanged.
  \node[bpmn event, label=above:{\scriptsize $p_{11}$}] (p11n) at (17.7,1.8) {};
  \node[bpmn event] (p8) at (12.9,0) {};
  \node[bpmn task, rounded corners=0pt, font=\tiny\sffamily, align=center] (csw) at (14.1,0) {check\\ sw};
  \node[bpmn event, label=below:{\scriptsize $p_9$}] (p9) at (15.3,0) {};
  \node[bpmn task, rounded corners=0pt, font=\tiny\sffamily, align=center] (cfg) at (16.5,0) {configure\\ sw};
  \node[bpmn event, label=below:{\scriptsize $p_{12}$}] (p12n) at (17.7,0) {};
  % Shared mutex place p10: holds one token so exactly one of the hw/sw
  % branches runs at a time: realises the ``sequentially, arbitrary order''
  % requirement of the specification (slides p.135--137).
  \node[bpmn event, label={[label distance=-1pt]below:{\scriptsize $p_{10}$}}] (p10m) at (15.3,0.9) {};
  \fill[accent] (p10m.center) circle (1.6pt);
  \node[bpmn task, rounded corners=0pt, font=\tiny\sffamily, align=center] (rd) at (18.9,0.9) {repair\\ done};
  \node[bpmn event] (p13) at (20.1,0.9) {};
  \node[bpmn task, rounded corners=0pt, font=\tiny\sffamily] (tst) at (21.2,0.9) {test};
  \node[bpmn event, label=above:{\scriptsize $p_{14}$}] (p14) at (22.3,0.9) {};
  \node[bpmn task, rounded corners=0pt, font=\tiny\sffamily, align=center] (ok) at (23.5,0) {test\\ ok};
  \node[bpmn task, rounded corners=0pt, font=\tiny\sffamily, align=center] (err) at (22.3,2.4) {error\\ detected};
  \node[bpmn event, label=below:{\scriptsize $o$}] (o) at (24.8,-1.3) {};
  \draw[bpmn flow] (i) -- (rc); \draw[bpmn flow] (rc) -- (p1);
  \draw[bpmn flow] (p1) -- (cd); \draw[bpmn flow] (cd) -- (p2);
  \draw[bpmn flow] (p2) -- (ec); \draw[bpmn flow] (ec) -- (p3);
  \draw[bpmn flow] (p3) -- (te); \draw[bpmn flow] (p3) -- (ac);
  \draw[bpmn flow] (te) -- (o);
  \draw[bpmn flow] (ac) -- (p5); \draw[bpmn flow] (p5) -- (sr);
  \draw[bpmn flow] (sr) -- (p6); \draw[bpmn flow] (sr) -- (p8); \draw[bpmn flow] (sr) -- (p10m);
  \draw[bpmn flow] (p6) -- (chw); \draw[bpmn flow] (chw) -- (p7); \draw[bpmn flow] (p7) -- (rhw); \draw[bpmn flow] (rhw) -- (p11n);
  \draw[bpmn flow] (p8) -- (csw); \draw[bpmn flow] (csw) -- (p9); \draw[bpmn flow] (p9) -- (cfg); \draw[bpmn flow] (cfg) -- (p12n);
  % mutex arcs: check-* borrow the p10 token, repair-*/configure-* return it
  \draw[bpmn flow] (p10m) -- (chw); \draw[bpmn flow] (p10m) -- (csw);
  \draw[bpmn flow] (rhw) -- (p10m); \draw[bpmn flow] (cfg) -- (p10m);
  \draw[bpmn flow] (p10m) -- (rd);
  \draw[bpmn flow] (p11n) -- (rd); \draw[bpmn flow] (p12n) -- (rd);
  \draw[bpmn flow] (rd) -- (p13); \draw[bpmn flow] (p13) -- (tst); \draw[bpmn flow] (tst) -- (p14);
  \draw[bpmn flow] (p14) -- (ok); \draw[bpmn flow] (ok) -- (o);
  \draw[bpmn flow] (p14) -- (err);
  \draw[bpmn flow] (err.west) -- (10.3,2.4) -- (p5);
\end{tikzpicture}}
\end{center}
The construction proceeds clause by clause: a linear intake (\emph{receive computer}, \emph{check defects}, \emph{estimate cost}); the customer's XOR at $p_3$ (\emph{too expensive} exits to $o$, \emph{acceptable cost} continues); the two repair activities; and the test loop: \emph{test ok} exits, \emph{error detected} returns to $p_5$ for another repair round. The specification is precise about the repair: the hardware and software activities run ``sequentially but in an arbitrary order'': not in parallel. A pure AND-split would let them interleave; instead \emph{start repair} marks all three of $p_6$, $p_8$ and the \textbf{shared place} $p_{10}$, and each \emph{check} consumes $p_{10}$ while the matching \emph{repair}/\emph{configure} returns it. That single token is a \textbf{mutex}: whichever branch starts first holds $p_{10}$ until its \emph{repair}/\emph{configure} completes, so the other cannot begin meanwhile. \emph{Repair done} then joins $p_{11}$, $p_{12}$ and $p_{10}$. The result is a WF net by construction (unique $i$, unique $o$, all nodes on $i$--$o$ paths). WoPeD confirms \textbf{soundness}: boundedness and liveness of $N^*$ hold, no dead or non-live transitions. The structural pane reports one free-choice violation and three PT-handles (with three TP-handles): not soundness defects, but patterns that block the deeper structural theorems (free-choice nets, S/T-components) of the next chapters.
\end{examplebox}
